%% =====================================================================
%%  T-01-03  Power as an Amoral Instrument
%%  -------------------------------------------------------------------
%%  One idea per file. Core is mandatory. Slots in canonical order.
%%  Max 700 words. No layout commands. No filler. New source material
%%  goes in sources/<BOOK>/T-01-03.tex -- never by rewriting this file.
%% =====================================================================
%% @kind: topic
%% @id: T-01-03
%% @title: Power as an Amoral Instrument
%% @subtitle: The mechanics do not care what you intend
%% @chapter: c01
%% @order: 30
%% @status: stable
%% @sources: B3
%% @updated: 2026-09-19

\topic{T-01-03}{Power as an Amoral Instrument}{The mechanics do not care what you intend}

\begin{core}
The strategic layer of this book treats power as a set of mechanics that operate whether or not you approve of them, and whether or not you intend to use them. Refusing to learn them does not remove them from the environment; it only removes them from your side.
\end{core}
\begin{mechanism}
The argument is descriptive rather than exhortatory. Court dynamics --- who is rising, who is exposed, who depends on whom --- exist in any hierarchy, and they continue to operate on a person who declines to notice them. What the practitioner adds is attention and choice: seeing the mechanism, and deciding deliberately whether to run it, run against it, or step out of its reach. The moral question is therefore not whether the mechanism exists but what you do once you can see it, which is a question this entry does not answer and the rest of the book keeps returning to.
\end{mechanism}
\begin{tells}
  \item You are surprised by an outcome that several other people saw coming.
  \item You describe an arrangement as unfair and nobody in it disagrees, and nothing changes.
  \item You believe that declining to play is a neutral position rather than a move with a payoff of its own.
\end{tells}
\begin{moves}
  \item Describe the situation in terms of interest and dependence rather than in terms of fairness.
  \item Identify who benefits from the current arrangement and what it costs them to change it.
  \item Only then decide whether to work within it, against it, or leave it.
\end{moves}
\begin{counters}
  \item Recognising the mechanism is itself the main defence: an arrangement you can describe in terms of interest is one you can decline, price, or exit.
  \item Refuse the inference that because something is amoral you must be. Seeing the game clearly and playing it badly on purpose are both available options.
  \item Keep at least one relationship or resource outside the system you are analysing, so that leaving is real rather than rhetorical.
\end{counters}
\begin{cost}
This framing has a known failure mode: it makes every interaction look strategic, which is both untrue and corrosive. Most exchanges are cooperative and most people are not running anything. Held as a permanent lens rather than a diagnostic tool, it produces isolation and misreading, and the book's own later material on isolation and on the cost of using people is a correction of it. It is also a description of how power behaves, not a licence; several of the moves described elsewhere in this handbook are crimes in most jurisdictions.
\end{cost}

\sources{B3}
\seesources
\seealso{T-01-01, T-01-04, T-27-01}
